\section{Conclusion}\label{sec:conclusion}

We presented \sys{}, a high-fidelity smart home simulation platform
that integrates Habitat~3.0 embodied AI simulation, QEMU-based
firmware-in-the-loop emulation, LLM-driven activity generation, and
bi-directional SmartThings cloud synchronization.

Our evaluation across five research questions and a large-scale
autonomous study (28~HSSD scenes, 7~simulated days each, 88.0~hours
wall-clock time) demonstrated that:
\begin{itemize}[leftmargin=*]
  \item \sys{}-generated activity schedules exhibit moderate
    distributional similarity to real-world CASAS and ARAS datasets
    (mean JS divergence~$= 0.218$),
    with the closest match on ARAS-House~B (JS~$= 0.101$).
    A leave-one-out Markov baseline comparison shows \sys{}
    outperforms Markov on diverse multi-occupant homes (ARAS),
    achieving 50--62\% lower JS divergence, while Markov excels
    on structurally similar single-occupant CASAS homes
    (\textbf{RQ1}).
  \item The platform scales to 200 concurrent event-bus devices at
    over 10{,}000~events/s with CPU below 36\% and memory under
    106\,MB, confirmed by 28~scenes of continuous operation
    with bounded memory growth (\textbf{RQ2}).
  \item Event-bus dispatch latency is 7\,$\mu$s at P99; database
    operations complete in under 3\,ms.  The autonomous evaluation
    delivered 47{,}207~SmartThings cloud updates with zero data
    loss (\textbf{RQ3}).
  \item Zero-shot Sim2Real transfer with enriched 119-dimensional
    feature representations achieves 0.813~accuracy on ARAS-A
    (F1~$= 0.783$), demonstrating that \sys{}-generated patterns
    capture transferable temporal structure.  Domain discrimination
    remains high on CASAS datasets due to categorical mismatch,
    motivating future domain adaptation work (\textbf{RQ4}).
  \item A security assessment of 982~attacks across five suites
    (firmware, network, phantom-delay, malicious SmartApp, ESP32
    buffer overflow) against six device types yields a 67.4\%
    exploit rate with a mean CVSS~3.1 score of~8.1, validated
    against 24~real-world CVEs and covering 7 of 10 OWASP IoT
    Top~10 categories.
    Packet-level traffic analysis---validated by \texttt{tshark}
    pcap captures (154{,}151~TCP segments, 12.1\,MB, 736~individual
    pcap files across 980~attack executions in 28~scenes,
    \Cref{sec:eval-pcap})---confirms that all attack traffic
    traverses the real Docker bridge network and is independently
    verifiable with Wireshark, substantiating \sys{}'s claim as a
    network-realistic security evaluation platform (\textbf{RQ5}).
\end{itemize}

The autonomous evaluation further validated end-to-end reliability:
94.9\% navigation success across 3{,}580~trials,
4{,}307~LLM-generated tasks using GPT-OSS~20B, and
26~of~28 scenes completing full navigation protocols without
intervention.  The two remaining scenes ran security attacks
successfully but failed to generate valid navmeshes.

\sys{} is, to our knowledge, the first platform to integrate embodied
3D simulation with real device firmware execution and live
cloud-connected IoT in a single reproducible framework.  A case study reproducing
the Phantom-Delay Attack~\cite{phantom-delay} within \sys{}'s
Docker network further validates the platform's utility for
IoT security research---enabling reproduction of published
protocol-level exploits without physical hardware.
We release the platform,
evaluation framework, and all experiment configurations as open
source to support reproducible IoT research.

\paragraph{Future Work.}
We plan to (1)~add multi-occupant coordination with inter-agent
communication, (2)~integrate additional IoT platforms (Google Home,
Apple HomeKit), (3)~conduct an ablation study across multiple LLM
models to quantify the impact of model choice on schedule diversity
and correctness, (4)~fine-tune a compact LLM on successful \sys{}
schedules to reduce generation errors, (5)~extend the firmware
library to include Zigbee and Z-Wave protocol stacks, and
(6)~investigate domain adaptation techniques (e.g., feature alignment,
few-shot fine-tuning) to further improve Sim2Real transfer accuracy
on fine-grained activity datasets beyond the 0.813 achieved on
ARAS-A.
